% Causal Probabilistic Transformer — preprint
%
% Format deliberately mirrors Wu & Tu, "Probabilistic Transformer: A Probabilistic Dependency
% Model for Contextual Word Representation" (Findings of ACL 2023) — same style file, same
% section skeleton, same length budget (8 pages of main body, everything else in appendices).
% See PAPER_STATUS.md for the section-by-section mapping and the page budget.
%
% Build: make   (output build/main.pdf)

\documentclass[11pt]{article}

% "preprint" = non-anonymous with page numbers. Switch to "review" for anonymous submission,
% "final" for camera-ready.
\usepackage[preprint]{acl}

\usepackage{times}
\usepackage{latexsym}
\usepackage[T1]{fontenc}
\usepackage[utf8]{inputenc}
\usepackage{microtype}
\usepackage{inconsolata}
\usepackage{graphicx}

% Packages and page setup. Kept deliberately venue-neutral: the document is a preprint
% first, and a conference style file can be dropped in later without touching the sections.

\usepackage[utf8]{inputenc}
\usepackage[T1]{fontenc}
\usepackage[margin=1in]{geometry}

\usepackage{amsmath}
\usepackage{amssymb}
\usepackage{amsthm}
\usepackage{mathtools}
\usepackage{bm}

\usepackage{graphicx}
\usepackage{booktabs}
\usepackage{multirow}

\usepackage{natbib}
\usepackage[colorlinks=true,allcolors=blue]{hyperref}
\usepackage[capitalise,noabbrev]{cleveref}

\usepackage{tikz}
\usetikzlibrary{arrows.meta,positioning,fit,calc}

% --- theorem environments ---
\theoremstyle{plain}
\newtheorem{theorem}{Theorem}
\newtheorem{proposition}[theorem]{Proposition}
\newtheorem{lemma}[theorem]{Lemma}
\newtheorem{corollary}[theorem]{Corollary}

\theoremstyle{definition}
\newtheorem{definition}[theorem]{Definition}
\newtheorem{remark}[theorem]{Remark}

% --- drafting aids -------------------------------------------------------------
% \TODO renders visibly in the draft. Set \draftfalse before circulating the PDF;
% every outstanding item then disappears, which is also a way to see how many remain.
\newif\ifdraft
\drafttrue

\ifdraft
  \usepackage{xcolor}
  \newcommand{\TODO}[1]{\textcolor{red}{\textbf{[TODO: #1]}}}
  \newcommand{\NOTE}[1]{\textcolor{blue}{\textbf{[note: #1]}}}
  % Marks a claim that must be checked against the reference PDFs before submission.
  \newcommand{\FROMPDF}[1]{\textcolor{orange}{\textbf{[verify against PDF: #1]}}}
\else
  \newcommand{\TODO}[1]{}
  \newcommand{\NOTE}[1]{}
  \newcommand{\FROMPDF}[1]{}
\fi

% Notation macros.
%
% The names follow the reference documents, NOT transformer convention. In particular
% `d` is the size of the label set of a word variable, not a hidden dimension, and `h`
% counts channels (head variables per word), not attention heads in the usual sense.
% Keep this file in sync with src/config.py, which uses the same names.
%
% Every macro below must correspond to something that exists in the factor graph — the
% modelling constraint of the project is that no parameter names a non-factor, and the
% notation should not be able to express one.

% --- index sets and sizes ---
\newcommand{\Vocab}{\mathcal{V}}                 % vocabulary
\newcommand{\nlab}{d}                            % number of labels of a word variable
\newcommand{\nchan}{h}                           % number of channels (head variables per word)
\newcommand{\krank}{r}                           % Kruskal rank of the arc-score decomposition
\newcommand{\nglob}{m}                           % number of global variables (Appendix B.3)
\newcommand{\clip}{\gamma}                       % distance clip threshold
\newcommand{\niter}{T}                           % MFVI iterations / loop depth
\newcommand{\seqlen}{n}                          % sequence length
\newcommand{\Dep}[1]{\mathcal{D}_{#1}}           % dependency candidate set at position #1

% --- random variables ---
\newcommand{\Wv}[1]{W_{#1}}                      % word variable (LATENT, not observed)
\newcommand{\Zv}[1]{Z_{#1}}                      % label variable
\newcommand{\Hv}[2]{H_{#1}^{(#2)}}               % head variable at position #1, channel #2
\newcommand{\Gv}[1]{G_{#1}}                      % global variable

% --- factors and parameters (each names a factor; see CLAUDE.md constraint 1) ---
\newcommand{\Sfac}{S}                            % word--label factor, |V| x d; mediates BOTH directions
\newcommand{\Tfac}[1]{T^{(#1)}}                  % arc score of channel #1, T = U V^\top
\newcommand{\Ufac}[1]{U^{(#1)}}
\newcommand{\Vfac}[1]{V^{(#1)}}
\newcommand{\bfac}{b}                            % word unary; equals the LM-head bias
\newcommand{\rfac}[1]{r^{(#1)}}                  % root key of channel #1
\newcommand{\Bfac}[1]{B^{(#1)}}                  % contracted arc--label score entering the readout

% --- variational quantities ---
\newcommand{\Q}{Q}                               % variational posterior
\newcommand{\QW}[1]{\Q_{\Wv{#1}}}
\newcommand{\QZ}[1]{\Q_{\Zv{#1}}}
\newcommand{\QH}[2]{\Q_{\Hv{#1}{#2}}}
\newcommand{\Free}{\mathcal{F}}                  % mean-field free energy
\newcommand{\lamZ}{\lambda_Z}                    % entropic Frank--Wolfe message weights
\newcommand{\lamH}{\lambda_H}
\newcommand{\lamW}{\lambda_W}
\newcommand{\lamG}{\lambda_G}

% --- readout ---
\newcommand{\readout}[1]{\mu_{#1}}               % \mu_t(a): unnormalised readout score of word a
\DeclareMathOperator*{\LSE}{LSE}                 % log-sum-exp
\DeclareMathOperator*{\logcumsumexp}{logcumsumexp}

% --- generic ---
\newcommand{\R}{\mathbb{R}}
\newcommand{\E}{\mathbb{E}}
\newcommand{\Prob}{\mathbb{P}}
\newcommand{\KL}[2]{\mathrm{KL}\!\left(#1 \,\|\, #2\right)}
\newcommand{\softmax}{\operatorname{softmax}}
\newcommand{\prefix}[1]{{<}#1}                   % "the prefix before position #1"
\newcommand{\stopgrad}[1]{\left\lfloor #1 \right\rfloor}  % frozen condition: no gradient path


\title{Causal Probabilistic Transformer:\\
       A Probabilistic Dependency Model for Autoregressive Language Modelling}
% NOTE: the title deliberately echoes Wu & Tu's, signalling that this is the causal counterpart
% of that model rather than an unrelated construction.

\author{
  Viktor Beliakov \\
  \And
  Penghao Kuang \\
  \And
  Kewei Tu \thanks{Corresponding author.} \\
}
% TODO(authors): confirm the author list, order and affiliation block with Prof. Tu before any
% public posting. Wu & Tu use a single shared-institution block (ShanghaiTech); match it.

\begin{document}
\maketitle

\begin{abstract}
% Wu & Tu's abstract is a single paragraph of about 150 words, and it ends on the "we hope this
% informs transformers" note rather than on a number. Match both. Written last; the version
% below fixes the claim so the rest of the paper has a target.

The probabilistic transformer of \citet{wu2023probabilistic} models contextual word
representation as posterior inference in a conditional random field over dependency structure,
and its mean-field inference procedure turns out to resemble a transformer encoder. That model
is an encoder only. In this work we derive its autoregressive counterpart. We show that
causality, a single global energy function and contextuality cannot hold simultaneously, and
resolve the conflict with a directed chain of conditional random fields in which the prefix
enters each step as a frozen condition. We further obtain a language-model head without adding
any parameter outside the factor graph, by promoting the word variable from observed to latent;
tied input and output embeddings then follow as a consequence rather than a design choice.
Because the per-position subgraph is a tree, the vocabulary readout is available in closed form.
Experiments on \TODO{corpora} compare our model against a standard causal transformer and
against a looped transformer, the latter isolating the effect of syntactic structure from that
of weight sharing. \TODO{one sentence of result} We hope that this work helps extend the
probabilistic account of contextual representation to generation.

\end{abstract}

\section{Introduction}
\label{sec:intro}

% Wu & Tu's introduction is about 1.2 columns and follows a fixed rhythm: (1) the field moved
% away from syntax; (2) transformers won; (3) they were designed heuristically; (4) we derive
% one instead; (5) what we find; (6) what we hope. Match the rhythm -- it is what makes the
% paper read as the sequel it is -- but replace the content with the causal question.

\TODO{draft. Written after the experiments, so the framing matches what was measured.}

\paragraph{Paragraph plan.}
\begin{enumerate}
  \item \citet{wu2023probabilistic} derive a transformer encoder rather than designing one: a
        CRF over dependency structure whose mean-field updates, unrolled, reproduce the
        computation graph of a transformer, with attention emerging as a posterior over
        syntactic heads rather than being posited.
  \item That model is an encoder. Autoregressive generation --- the setting in which
        transformers are actually used --- has no counterpart in the framework, and
        \citet{wu2023probabilistic} name a decoder as an open extension.
  \item Masking the factors does not produce one. \Cref{sec:obstruction} states why: causality,
        a single global energy and contextuality cannot hold together.
  \item Our model: a directed chain of conditional CRFs in which the prefix enters as a frozen
        condition, so the anti-causal term is never formed rather than removed afterwards.
  \item The output mechanism, which is the part to check: the language-model head introduces
        \emph{no} parameter. The word variable, already in the graph, is simply unclamped.
  \item What we find empirically, and the framing --- structure versus weight sharing, not
        versus perplexity.
\end{enumerate}

\paragraph{One principle, two results.}
\NOTE{This paragraph is the one that keeps the paper from reading as two loosely related
contributions. Without it, \Cref{sec:obstruction} and \Cref{sec:output} are an impossibility
theorem and an engineering trick that happen to share a paper.}
Both of our main results are refusals, and they refuse for the same reason: the factor graph is
the model, and it constrains what may be added to it. \Cref{sec:obstruction} says one may not
keep a single global energy and still be causal. \Cref{sec:output} says one may not add an
output matrix that names no factor. Taking the graph seriously in the first case forces the
chain of conditional CRFs; taking it seriously in the second forces tied embeddings, an output
bias equal to the word unary term, and the identity of the encoder and the decoder as two
conditioning patterns of one model.

\paragraph{Contributions.}
\begin{itemize}
  \item A statement of what a causal probabilistic transformer cannot be
        (\Cref{sec:obstruction}), and a construction satisfying the remaining constraints
        (\Cref{sec:chain}).
  \item A graph-faithful output mechanism (\Cref{sec:output}) whose consequences --- forced
        embedding tying, the output bias, one model for both directions --- are derived rather
        than chosen.
  \item An exact readout (\Cref{prop:exact-readout}): the per-position subgraph is a tree, so
        sum-product is exact, and along the sequence it is a single causal
        \texttt{logcumsumexp} scan that preserves the layer-parallel schedule.
  \item An empirical study (\Cref{sec:experiments}) that separates the contribution of
        structure from that of weight sharing by comparing against a Looped transformer as well
        as a standard decoder.
\end{itemize}

\paragraph{What this paper does not claim.}
We do not claim to improve on a standard decoder in perplexity, and the experiments are not
designed to. \citet{wu2023probabilistic} state that the aim of the model is to inform
transformers rather than compete with them, and report degradation beyond roughly $10^5$
sentences; we inherit both the framing and the regime. There is in addition an honest
rank-$\nlab$ bottleneck: all predictive information passes through the label variables. The
question we ask is whether syntactic structure carries information that weight sharing alone
does not, which is a question about the gap to a Looped transformer rather than the gap to a
standard decoder.

\section{Causal Probabilistic Transformers}
\label{sec:model}

% Mirrors Wu & Tu Section 2 ("Probabilistic Transformers"): the model first, then inference,
% then variants. Their Section 2 runs about 2.3 pages of a two-column ACL page; budget the
% same here. Figure 1 (the factor graph) belongs in 2.2, early, exactly as theirs does.

We first recall the model of \citet{wu2023probabilistic} only as far as we need it
(\Cref{sec:recap}), then state the obstruction that prevents a direct causal version
(\Cref{sec:obstruction}), then give the model (\Cref{sec:chain}), its output mechanism
(\Cref{sec:output}), its inference procedure (\Cref{sec:inference}), and finally some variants.

\subsection{Background}
\label{sec:recap}

\FROMPDF{Wu \& Tu \S2.1--\S2.2 --- variable set, the two potential functions, MFVI}
\TODO{compress to about half a column. Only what \Cref{sec:obstruction,sec:output} actually
use: the label variables $\Zv{i}$ over $\nlab$ labels; the head variables $\Hv{i}{c}$ over
candidate heads, one set per channel; the unary factor $\phi_u(\Zv{i}) = \exp(\Sfac_{w_i,\Zv{i}})$
with $\Sfac \in \R^{|\Vocab| \times \nlab}$; the ternary factor over $(\Hv{i}{c}, \Zv{i}, \Zv{j})$
with score matrix $\Tfac{c}$; and that MFVI unrolled for $\niter$ iterations has the shape of a
$\niter$-layer transformer encoder with shared weights. Note explicitly that the word $w_i$
enters \emph{only} through $\Sfac$ --- \Cref{sec:output} turns on that fact.}

\subsection{Why a causal mask is not enough}
\label{sec:obstruction}

\TODO{about half a column. Wu \& Tu's paper contains no theorem environments; keep this to one
theorem statement in prose-adjacent form and push the proof to \Cref{app:proof}.}

The obvious route to a causal model is to mask the ternary factors so that $\Hv{i}{c}$ may only
point to positions $j \le i$. This does not give a causal model. Informally, three properties
that one would want cannot hold together:

\begin{itemize}
  \item \textbf{Causality:} the posterior at position $i$ depends only on $w_{1:i}$;
  \item \textbf{A single global energy:} one energy function over the sentence, normalised once;
  \item \textbf{Contextuality:} the representation at $i$ genuinely depends on its prefix.
\end{itemize}

\begin{theorem}[Trilemma]
\label{thm:trilemma}
\TODO{formal statement} No model satisfies all three simultaneously.
\end{theorem}

The mechanism is the partition function: a factor coupling positions $i<j$ contributes to the
marginal at $i$ whenever it contributes to the marginal at $j$, and forbidding the first while
keeping the second is inconsistent with a single normalisation. Masking the factors changes
which configurations have non-zero score; it does not stop the surviving ones from being
normalised jointly. The proof is in \Cref{app:proof}.

Causality and contextuality are what the model is for, so the single global energy is what we
give up.

\begin{remark}[Never formed, not deleted]
\label{rmk:never-formed}
Building the symmetric model and then removing the anti-causal contribution gives a different
object from building a model in which that contribution never exists. This distinction is what
\Cref{sec:chain} implements and what \Cref{sec:exp-correctness} tests.
\end{remark}

\subsection{A directed chain of conditional CRFs}
\label{sec:chain}

\FROMPDF{Part~II of the main document --- the factorisation and the conditioning statement}

\begin{figure}[t]
\centering
\TODO{Figure 1. Two panels of the \emph{same} factor graph: (a) encoder, $\Wv{i}$ shaded
(observed); (b) decoder, $\Wv{i}$ unshaded (latent), with the prefix boxed as a frozen
condition. Same factors, same edges, different shading. This single figure carries the claim
that encoding and decoding are two conditioning patterns of one model, and it is the analogue
of Wu \& Tu's Figure 1, which is also introduced at exactly this point.}
\caption{The factor graph of the causal probabilistic transformer for $\seqlen = 3$.}
\label{fig:factor-graph}
\end{figure}

Instead of one CRF over the sentence, we define a chain of conditional CRFs, one per position.
The CRF at position $t$ ranges over $(\Zv{t}, \Hv{t}{1..\nchan}, \Wv{t})$ and is conditioned on
the prefix posteriors $\{\QZ{s}\}_{s<t}$, which enter as fixed numbers:
%
\begin{equation}
\label{eq:chain}
  \TODO{the factorisation}
\end{equation}

\paragraph{The prefix is a frozen condition.}
Each $\QZ{s}$, $s<t$, is data at step $t$: no message travels from $t$ back to $s$. The
anti-causal term of \Cref{thm:trilemma} is therefore never formed, and causality holds by
construction rather than by masking.

\paragraph{Frozen forward, not frozen backward.}
\emph{Frozen} means constant with respect to step $t$'s inference problem, not detached in
autodiff. Training gradients flow backwards through the whole prefix computation, exactly as
they do through cached activations in a causal transformer; this is how $\Sfac$ is trained in
both of its roles (\Cref{sec:output}). \NOTE{say this explicitly in the paper. It is the single
easiest thing for a reader to get backwards, and getting it backwards makes the model look
either non-causal or untrainable.}

\subsection{The output mechanism}
\label{sec:output}

This is the part of the construction we most want checked, and we present it as such.

\paragraph{The design we rejected.}
The obvious way to obtain next-token logits is to project $\QZ{t}$ through a fresh matrix in
$\R^{\nlab \times |\Vocab|}$. Such a matrix corresponds to no factor in the graph: the output
would not traverse the factor graph at all, and the model would be a CRF with a neural head
attached. We report this because rejecting it is what forced the mechanism below.

\paragraph{Unclamping the word variable.}
The word variable is already in the graph and $\Sfac$ is already a binary factor; in the encoder
it is evaluated with its word argument clamped to the observed token. The output mechanism is to
stop clamping it:
%
\begin{equation}
\label{eq:unclamp}
  \TODO{$\readout{t}(a) \propto \dots$ in the notation of \texttt{causal\_pt\_output\_note.pdf}}
\end{equation}
%
No factor is added and no parameter is created. Next-token prediction is the posterior over a
variable the model already had.

\paragraph{Consequences.}
Three properties follow, and we stress that they are consequences rather than design choices:

\begin{enumerate}
  \item \textbf{Tying is forced.} One factor $\Sfac$ mediates both directions, so input and
        output word embeddings are the same matrix. Untying is not a hyperparameter; it is a
        different model.
  \item \textbf{The output bias is the word unary.} \TODO{state}
  \item \textbf{One model, two conditioning patterns.} The encoder clamps $\Wv{i}$; the decoder
        unclamps it and directs the chain. Same parameters, same factors
        (\Cref{fig:factor-graph}).
\end{enumerate}

\subsection{Inference}
\label{sec:inference}

\NOTE{Resolved 2026-08-05: \S17.2 proposes MFVI as the mainline but \S23.3 supersedes it. The
exact readout is the model; the mean-field readout is the ablation of \Cref{sec:exp3}. Do not
present them as two equal options.}

\FROMPDF{Part~III \S16--\S17.1 --- update equations and schedule}
\begin{align}
\label{eq:update-Z}
  \QZ{t} &\propto \TODO{} \\
\label{eq:update-H}
  \QH{t}{c} &\propto \TODO{}
\end{align}

\paragraph{Exact readout.}
The subgraph at position $t$ is a star centred on $\Zv{t}$, hence a tree, so sum-product is
exact: no approximation and no iteration are involved.

\begin{proposition}
\label{prop:exact-readout}
\FROMPDF{\S23.3 --- confirm the form and the seeding by $\rfac{c}$}
\begin{equation}
\label{eq:exact-readout}
  \log \readout{t}(a) = \sum_{c=1}^{\nchan} \LSE_{j \in \Dep{t}} \Bfac{c}_{j,a},
\end{equation}
seeded with the root key $\rfac{c}_a$.
\end{proposition}

Along the sequence $\Dep{t}$ grows by one position per step, so \Cref{eq:exact-readout} is a
causal prefix log-sum-exp: one \texttt{logcumsumexp} scan, $O(\seqlen\nlab\nchan)$, fully
parallel over positions. The exact readout therefore costs neither the layer-parallel schedule
nor an inner loop --- with it there is no query stream and no predictive iteration at all. Two
costs are real and we state them rather than bury them: the $|\Vocab| \times \nlab$ readout is a
log-sum-exp rather than a matmul (same FLOPs, worse hardware constants, chunked over the
vocabulary with a fused cross-entropy), and LSE-pooling gradients are untested at
language-model scale.

\paragraph{Every step is valid inference.}
\FROMPDF{Part~III \S18 Check~3, as restated in \S23.3}
Mean field where exactness is impossible (the chain), exact sum-product where the subgraph is a
tree (the slot).

\subsection{Extensions and variants}
\label{sec:variants}

% Wu & Tu keep this subsection short and push the rest to their Appendix B. Do the same.

\paragraph{Distance.} \TODO{the causal half of the clipped ternary potential: only $i-j>0$
exists, so there are $\clip+1$ buckets rather than $2\clip+1$.}

\paragraph{Update schedule.} \TODO{parallel (a depth-$\niter$ shared causal transformer) vs.\
serial ($\tau$ inner rounds per slot); say which produced the reported numbers.}

\paragraph{Global variables.} \TODO{the Appendix B.3 globals $\Gv{t}$, proposed by
\citet{wu2023probabilistic} as an in-graph substitute for the missing feed-forward structure and
never tested there. We treat them as a measured variable (\Cref{sec:exp-ablations}), not an
assumption.}

Further variants are in \Cref{app:variants}.

\section{Comparison with Causal Transformers}
\label{sec:comparison}

% This section is the analogue of Wu & Tu Section 3 ("Comparison with Transformers"), which is
% the LARGEST section of their paper -- about 2.4 pages, as long as their experiments. It had
% no counterpart in the first outline of this paper, and that was the biggest structural gap:
% in this format the conceptual comparison carries as much weight as the numbers, and it is
% also where the Looped baseline of Section 4 gets motivated rather than asserted.

\citet{wu2023probabilistic} show that MFVI in the encoder model reproduces the computation graph
of a transformer encoder, with the posterior over dependency heads playing the role of
self-attention. We ask the same question of the causal model, and the answer is the reason the
experiments in \Cref{sec:experiments} take the shape they do.

\subsection{Tensor form}
\label{sec:tensor-form}

\TODO{re-derive the update equations of \Cref{sec:inference} in tensor form, following the
structure of Wu \& Tu \S3.1 so the two papers can be read side by side. Their $Q_z^{(t)}$,
$Q_{h,c}^{(t)}$, $\mathcal{F}$, $\mathcal{G}$ notation should be reused verbatim wherever the
quantity is the same one --- deviating from it without reason costs the reader the comparison.}

\subsection{Single-channel update vs.\ masked self-attention}
\label{sec:vs-attention}

\TODO{the causal analogue of their \S3.2. The claim to establish: the single-channel update
under the causal chain reduces to scaled dot-product attention with a causal mask, and the
correspondence $Q_c \leftrightarrow W^Q$, $K_c = V_c \leftrightarrow W^K = W^V$ carries over.
State precisely where it stops holding --- that is more interesting than where it does.}

\subsection{The readout vs.\ a language-model head}
\label{sec:vs-lm-head}

\TODO{no counterpart in Wu \& Tu, because the encoder has no LM head. Show that
\Cref{eq:exact-readout} is a mixture of $\nchan$ softmaxes over a $\nlab$-dimensional space,
compare with the ordinary $\R^{\nlab\times|\Vocab|}$ head, and connect to
\citet{yang2018breaking}: the rank-$\nlab$ bottleneck is real, but a mixture of softmaxes is not
rank-limited in the way a single softmax is. Whether that helps is measured in
\Cref{sec:exp3}.}

\subsection{Full model comparison}
\label{sec:full-comparison}

\begin{figure*}[t]
\centering
\TODO{Figure 2, full width. Three computation graphs side by side --- (a) GPT-style decoder,
(b) Looped decoder, (c) causal probabilistic transformer --- in the visual style of Wu \& Tu's
Figure 3. This figure is what makes \Cref{tab:three-point} unnecessary to argue for in words.}
\caption{Computation graphs for a causal transformer, a Looped transformer, and a causal
probabilistic transformer.}
\label{fig:computation-graphs}
\end{figure*}

\TODO{prose. The differences to draw out, following their \S3.4:}

\begin{itemize}
  \item \textbf{No feed-forward structure.} Inherited from the encoder model, and the reason
        the global variables of \Cref{sec:variants} are worth measuring.
  \item \textbf{No residual connections or layer norm.} The unary scores are re-added at each
        iteration instead.
  \item \textbf{Parameters shared across iterations.} This is the property the Looped
        transformer \citep{dehghani2019universal, giannou2023looped} also has, and it is why a
        comparison against a standard decoder alone cannot attribute any difference to
        structure.
  \item \textbf{The output is a posterior, not a projection.} No parameter exists outside the
        factor graph (\Cref{sec:output}).
\end{itemize}

\begin{table}[t]
\centering
\small
\begin{tabular}{lcc}
\toprule
Model & \makecell{Weight\\sharing} & \makecell{Syntactic\\structure} \\
\midrule
Causal transformer & no  & no  \\
Looped transformer & yes & no  \\
Causal PT          & yes & yes \\
\bottomrule
\end{tabular}
\caption{The three models of \Cref{sec:experiments}. Looped vs.\ transformer isolates weight
sharing; causal PT vs.\ Looped isolates structure, which is the comparison this paper is about.}
\label{tab:three-point}
\end{table}

\Cref{tab:three-point} is what the experimental design follows from: the causal PT differs from
a standard decoder in two ways at once, so a two-model comparison cannot attribute a difference
to either. \NOTE{putting this table in Section 3 rather than Section 4 is deliberate --- it
makes the Looped baseline a consequence of the model comparison instead of an experimental
choice that a reviewer could read as cherry-picked.}

\section{Experiments}
\label{sec:experiments}

% Mirrors Wu & Tu Section 4: three subsections (Tasks and Datasets / Settings / Results), about
% 1.6 pages, and -- the part worth copying -- ONE results table for the whole paper. Their
% Table 1 covers five tasks and six datasets in eleven rows, and their Results prose is barely
% half a column. Detailed dataset descriptions go to the appendix.

\NOTE{No number appears here without a row in the run log of the corresponding
\texttt{experiments/*/EXPERIMENT\_STATUS.md}: date, commit, config, seed, metric, wall-clock.}

\subsection{Tasks and Datasets}
\label{sec:tasks}

\TODO{brief, with the detail in \Cref{app:datasets} --- this is how Wu \& Tu handle it.}

\paragraph{Language modelling.} \TODO{PTB \citep{marcus1993building} and/or WikiText-2
\citep{merity2017pointer}; perplexity on the held-out set.}

We deliberately do not use WikiText-103. \citet{wu2023probabilistic} report that the encoder
model significantly underperforms transformers beyond roughly $10^5$ sentences, suspected to be
the absence of a feed-forward structure; running there would test a known failure mode of the
encoder rather than the causal extension. \NOTE{state this before a reviewer asks why the
corpus is small --- their own paper makes the same point in its Discussion}

\subsection{Settings}
\label{sec:settings}

\TODO{parameter budget 20--50M matched across models; identical tokenizer, vocabulary and
context length; one shared training loop; Adam with identical hyperparameters and schedule;
$k$ seeds, mean $\pm$ std. Wu \& Tu report mean and std over 5 random runs --- match that.}

\paragraph{Matching budgets.}
We report parameter count \emph{and} wall-clock, since the causal PT shares parameters across
iterations and equal parameter count therefore does not mean equal compute.

\paragraph{The embedding split.}
We report embedding and non-embedding parameters separately, which matters more here than in a
usual comparison. With tying forced (\Cref{sec:output}) the causal PT's parameters sit almost
entirely in $\Sfac \in \R^{|\Vocab|\times\nlab}$; the arc scores $\Tfac{c}$, the word unary
$\bfac$ and the root key $\rfac{c}$ are negligible beside it. At a matched total the causal PT
is close to an embedding table plus a little and a transformer baseline is not, and a single
total hides that completely. \TODO{state the matching rule chosen --- total parameters,
non-embedding parameters, or both with vocabulary fixed. Any is defensible; leaving it unstated
is what gets attacked first.}

\paragraph{Loop depth.}
The loop depth of the Looped baseline matches the number of MFVI iterations $\niter$, so
effective depth is comparable as well as parameter count. We report a sweep over $\niter$ rather
than a single value, so the comparison does not rest on one lucky setting.

\subsection{Results}
\label{sec:results}

\begin{table*}[t]
\centering
\small
\begin{tabular}{llcccc}
\toprule
Dataset & Metric & \makecell{Params\\(emb.\ / non-emb.)} & Transformer & Looped & Causal PT \\
\midrule
PTB         & Perplexity & \TODO{} & \TODO{} & \TODO{} & \TODO{} \\
WikiText-2  & Perplexity & \TODO{} & \TODO{} & \TODO{} & \TODO{} \\
\bottomrule
\end{tabular}
\caption{Main results, mean $\pm$ std over \TODO{$k$} random runs. The comparison this paper is
about is the last two columns.}
\label{tab:main}
\end{table*}

\paragraph{Correctness of the implementation.}
\label{sec:exp-correctness}
\TODO{two sentences plus \Cref{app:validation}. Only the two checks that test the derivation
rather than the code belong in the main text: the mean-field free energy is non-increasing
across iterations on a toy graph, and \Cref{eq:exact-readout} agrees with brute-force
enumeration over a tiny vocabulary to \TODO{1e-N}. Causality, tying, normalisation and the
single-batch overfit are software checks --- list them in the appendix.}

\paragraph{Language modelling.}
\TODO{read \Cref{tab:main}. Two readings, in this order: (i) the causal PT trains and lands
within a modest gap of the transformer --- necessary but not the point; (ii) the causal PT
vs.\ Looped gap is the paper's claim, because Looped has the weight sharing without the
structure (\Cref{tab:three-point}). If the causal PT is worse than Looped across the $\niter$
sweep, the central claim fails and we say so here.}

\paragraph{Ablations.}
\label{sec:exp-ablations}
\TODO{(i) global variables $\Gv{t}$ on/off --- a measured variable, not an assumption, and the
delta is itself a result since \citet{wu2023probabilistic} propose them but never test them;
(ii) exact vs.\ mean-field readout, below.}

\paragraph{Exact vs.\ mean-field readout.}
\label{sec:exp3}
\TODO{Two variants. (a) Evaluation-time swap: train with the mean-field readout, substitute
\Cref{eq:exact-readout} at evaluation on the same parameters --- isolates the approximation
error, nearly free. (b) Train with each --- measures the end-to-end effect, one extra full run.
Neither outcome is a win or a loss: a small gap says mean field is adequate and cheap, a large
one says the tree structure should be exploited. Also check whether the mixture-of-softmaxes
form relieves the rank-$\nlab$ bottleneck of \Cref{sec:vs-lm-head}.}

\section{Related Work}
\label{sec:related}

% Wu & Tu's Related Work is one third of a column -- a single dense paragraph. Match that
% length; the space belongs to Sections 2 and 3.

\TODO{one paragraph, in this order:}

\begin{itemize}
  \item \textbf{Syntax in transformers.} \TODO{the line \citet{wu2023probabilistic} cite ---
        supervised attention to syntactic governors, constituency/dependency injection,
        dependency-constrained attention. Position this work as the same deviation they claim:
        we start from probabilistic modelling and arrive at a transformer, rather than adding
        syntax to one.}
  \item \textbf{Weight sharing and looped models.} \TODO{\citet{dehghani2019universal},
        \citet{giannou2023looped}, \citet{yang2024looped}. This citation is load-bearing: it
        establishes that the Looped baseline of \Cref{sec:experiments} is a recognised model and
        not a strawman constructed for this paper.}
  \item \textbf{Structured autoregressive models.} \TODO{what exists for causal LMs with latent
        structure, and why none of it is a causal PT.}
  \item \textbf{The softmax bottleneck.} \TODO{\citet{yang2018breaking}; connect to
        \Cref{sec:vs-lm-head}.}
\end{itemize}

\section{Discussion}
\label{sec:discussion}

% Wu & Tu have a Discussion section separate from their Conclusion, and it is where the honest
% statements live: that the goal is not to compete with transformers, and that the model
% underperforms beyond ~100k sentences. We inherit both the framing and the section. This is
% the most useful structural thing copied from their paper -- it gives the concessions a place
% of their own, so they read as scope rather than as apology buried in a results paragraph.

\paragraph{What this paper does and does not claim.}
As in \citet{wu2023probabilistic}, our goal is not to propose a model that competes with
transformers on perplexity. It is to establish that the probabilistic-dependency account of
contextual representation extends to autoregressive generation, and to measure what its
structure contributes. \TODO{one sentence stating what the numbers actually showed, written
after they exist and not before.}

\paragraph{The rank-$\nlab$ bottleneck.}
\TODO{All predictive information reaches the vocabulary through $\Zv{t}$, whose support has size
$\nlab$. This is an honest cost of the construction rather than a defect of the implementation,
and it predicts a gap to a standard decoder that no amount of tuning removes. Say by how much,
and relate to the measurement in \Cref{sec:exp3}.}

\paragraph{The missing feed-forward structure.}
\TODO{inherited from the encoder model, and the leading suspect in its degradation on larger
corpora. Report what the global-variable ablation of \Cref{sec:exp-ablations} found --- this is
the first empirical test of the Appendix B.3 proposal of \citet{wu2023probabilistic}, and it is
worth saying so explicitly whichever way it came out.}

\paragraph{Extensions.}
\TODO{keep short. Candidates: continuous or vector-valued $\Zv{}$; dependency-label variables
paired with each $\Hv{}{}$; whether the exact readout suggests anything transferable back to
ordinary decoders --- which is the direction \citet{wu2023probabilistic} name as the point of
the whole programme.}

\section{Conclusion}
\label{sec:conclusion}

% Wu & Tu's Conclusion is one short paragraph that restates the model and the headline result,
% with no new material. Match that -- roughly 0.2 of a column.

\TODO{one paragraph, written last and written to the measured result. If the causal PT matches
Looped, the claim is that syntactic structure carries information weight sharing does not. If it
does not, the contribution is the construction plus a negative empirical finding, which is still
a paper --- do not reframe the numbers to rescue the stronger claim.}

% Unnumbered, placed after the Conclusion and before the references -- this is mandatory in the
% *ACL format since 2023, and it does not count against the page limit. Wu & Tu use it for three
% concrete concessions rather than generalities; do the same.

\section*{Limitations}
\label{sec:limitations}

\TODO{three or four concrete items. Candidates, each a consequence of the construction rather
than something more tuning would fix:}

\begin{itemize}
  \item \textbf{The rank-$\nlab$ bottleneck.} All predictive information reaches the vocabulary
        through the label variables (\Cref{sec:vs-lm-head}).
  \item \textbf{Scale.} Evaluated only on small corpora, for the reason given in
        \Cref{sec:tasks}; behaviour beyond that regime is untested here, and the encoder model
        is reported to degrade there.
  \item \textbf{No single sequence-level joint.} Giving up the global energy
        (\Cref{sec:obstruction}) means quantities defined by that joint are unavailable.
  \item \textbf{Speed.} \TODO{measure. The vocabulary readout is a log-sum-exp rather than a
        matmul: same FLOPs, worse constants. Wu \& Tu report their encoder as about 3 times
        slower than a transformer despite having far fewer parameters --- report the analogous
        number here rather than leaving the reader to assume it is small.}
  \item \TODO{at least one item that the runs exposed rather than the construction predicted.
        A limitations section made entirely of things known in advance is less credible than one
        that admits a surprise.}
\end{itemize}


\bibliography{refs}

\appendix
\section{Derivations}
\label{app:derivations}

% Mirrors Wu & Tu Appendix A, which carries the inference details the main text compresses
% (their A.5 is where the message weights \lambda_Z, \lambda_H are justified).

\subsection{Full mean-field updates}
\label{app:full-updates}
\TODO{the complete updates from Part~III \S16, with the steps \Cref{sec:inference} compresses.}

\subsection{Derivation of the exact readout}
\label{app:exact-derivation}
\TODO{sum-product on the star subgraph, giving \Cref{eq:exact-readout}; and the
\texttt{logcumsumexp} form along the sequence.}

\subsection{Message weights}
\label{app:message-weights}
\TODO{$\lamZ, \lamH, \lamW, \lamG$ --- follow Wu \& Tu \S A.5, and state where the causal model
forces a different choice from theirs.}

\subsection{Complexity}
\label{app:complexity}
\TODO{per-token cost of the content-stream solve, the $O(\nlab\nchan)$ cache update and the
vocabulary readout; and the chunked fused cross-entropy used in the implementation.}

\section{Additional Variants}
\label{app:variants}

% Mirrors Wu & Tu Appendix B, where the variants not used in the main experiments live --
% including the B.3 global variables that Section 2.6 promotes to a measured ablation here.

\subsection{Global variables}
\label{app:globals}
\TODO{the $\Gv{t}$ construction, following Appendix B.3 of \citet{wu2023probabilistic}, and
what changes in the causal setting. The measurement is in \Cref{sec:exp-ablations}; this is the
model description.}

\subsection{Serial update schedule}
\label{app:serial}
\TODO{$\tau$ inner rounds per observed slot; when it differs from the parallel schedule and by
how much.}

\subsection{Root node}
\label{app:root}
\TODO{whether the dummy root of Wu \& Tu \S2.3.5 survives in a causal model, and what the root
key $\rfac{c}$ of \Cref{eq:exact-readout} is in relation to it.}

\section{Proof of \Cref{thm:trilemma}}
\label{app:proof}

\TODO{the full statement and proof, from Part~I of the main document. \Cref{sec:obstruction}
carries only the statement and the mechanism in one paragraph, which is what the format allows;
everything else is here.}

\TODO{state the three properties formally as they are compressed to bullets in the main text:
causality, a single global energy, contextuality. Then the incompatibility argument.}

\section{Datasets}
\label{app:datasets}

% Wu & Tu's Appendix D. Their main text gives one sentence per dataset and defers everything
% else here; \Cref{sec:tasks} does the same.

\TODO{corpus statistics, preprocessing, tokenisation, vocabulary construction, train/dev/test
splits, and the exact sequence length. Enough for someone to reproduce the numbers in
\Cref{tab:main} without reading the code.}

\section{Implementation Validation}
\label{app:validation}

% No counterpart in Wu & Tu. It exists because two of these checks test the derivation rather
% than the code, and because a causal model that is silently non-causal produces entirely
% plausible numbers.

\TODO{the full checklist, with pass/fail and the numbers. \Cref{sec:exp-correctness} reports
only the two that test the equations; these are the rest.}

\begin{itemize}
  \item \textbf{Causality.} Changing the token at position $t$ leaves the logits at every
        position $<t$ bitwise unchanged (CPU, fixed seed --- on GPU, non-deterministic
        reductions can perturb the last bit of correct code).
  \item \textbf{Frozen prefix.} Prefix quantities are kept as attached leaves with
        \texttt{requires\_grad=True}; after backpropagating the loss at step $t$, their gradient
        is exactly zero for $j \ge t$ and non-zero for $j < t$. Testing this on a detached
        tensor would pass unconditionally and is therefore no test at all.
  \item \textbf{Tying.} Input and output embeddings are the same tensor object, not merely
        equal.
  \item \textbf{Normalisation.} All posteriors sum to one along their variable axis at every
        iteration.
  \item \textbf{Free energy.} The mean-field free energy is non-increasing across iterations on
        a toy graph. \TODO{figure} \NOTE{this is the check that catches a typo in the update
        equations; shapes, normalisation and causality are all blind to one.}
  \item \textbf{Exact readout vs.\ brute force.} On one slot the graph is a tree, so
        \Cref{eq:exact-readout} must agree with explicit enumeration over a tiny vocabulary.
        \TODO{agreement to 1e-N} This validates the factorisation itself.
  \item \textbf{Single-batch overfit.} Loss on one batch falls to near zero.
\end{itemize}

\section{Worked Example}
\label{app:worked-example}

% Wu & Tu's Appendix F holds case studies of the dependency structures their model infers. Ours
% holds the numeric example, which serves the same purpose: showing the reader what the model
% actually computes rather than asserting that it computes it.

\TODO{reproduce \texttt{causal\_pt\_output\_note.pdf} \S5 --- $\Vocab = \{\textit{the},
\textit{cat}, \textit{sat}, \textit{mat}\}$, $\nlab = 2$, one channel, both readout modes, with
printed posteriors and logits at every step.}

\NOTE{take the numbers from the test-suite output rather than retyping them, so the appendix and
the tests cannot drift apart, and say in the caption that they agree with the independently
derived reference.}

\TODO{if space allows, add the analogue of Wu \& Tu's Appendix F proper: a case study of the
dependency structures the causal model infers on real text. Their finding was that the induced
structures are only partially consistent with human intuition --- reporting the same honestly is
better than omitting the section.}


\end{document}
